\section{Related Work}

\xhdr{Time Series Reasoning}
Motivated by the need to solve real-world questions that require reasoning over evolving signals, as well as by advances in multimodal LLMs capable of processing non-textual inputs, recent studies have begun to formalize time-series reasoning. Several benchmarks have been introduced to evaluate such capabilities by pairing time-series signals with natural-language queries, including ECG-QA~\cite{ecgqa_dataset}, TSQA~\cite{tsqa_dataset}, TRQA~\cite{trqa_dataset}, RCW~\cite{rcw_dataset}, and related datasets~\cite{merrill2024tsandlang}. These resources frame temporal reasoning as question answering, explanation, or diagnostic inference over real or synthetic data, providing standardized settings for assessing LLM-based approaches.
In parallel with benchmark development, several modeling strategies have emerged to support time-series reasoning in LLMs. One line of work treats time-series values as textual inputs, allowing standard text-only LLMs to process temporal data without architectural modification. A second line couples a time-series encoder with a language model, such as ChatTS~\cite{chatts}, OpenTSLM~\cite{langer2025opentslm}, and ITFormer~\cite{itformer}, mapping numerical observations into learned embeddings that the LLM can condition on for reasoning. A third line of work encodes time series into visual representations. VL-Time~\cite{rcw_dataset} and TimeMaster~\cite{zhang2025timemaster} render signals as plots and use vision-language models to interpret them. These methods improve performance on time-series reasoning tasks, but they mainly address representation and modality integration, that is, how temporal signals are encoded and adapted for LLM processing. In contrast, the problem of selecting and analyzing temporal information during reasoning remains underexplored.

\xhdr{Dynamic Visual Search in Vision-Language Models}
A related line of work in vision-language reasoning focuses on developing models that iteratively select or zoom in on informative image regions conditioned on intermediate reasoning steps, acquiring additional visual information as needed to solve a task. Prompt-based methods such as ZoomEye~\cite{shen2025zoomeyeenhancingmultimodalllms} and Chain-of Spot~\cite{liu2024chainofspotinteractivereasoningimproves} use tree-based search or question-conditioned region-of-interest cropping to localize question-relevant regions, while Set-of-Mark prompting~\cite{yang2023setofmarkpromptingunleashesextraordinary} and DetToolChain~\cite{wu2024dettoolchainnewpromptingparadigm} leverage segmentation-based visual marks or detection-oriented prompting toolkits to ground reasoning in specific regions. More recent reinforcement learning approaches extend this paradigm with learned selection policies: Chain-of-Focus~\cite{zhang2025chain} trains a model to progressively zoom into regions based on intermediate cues, DeepEyes~\cite{zheng2026deepeyesincentivizingthinkingimages} interleaves multimodal chain-of-thought with native visual grounding, and Pixel-Reasoner~\cite{wang2025pixel} uses curiosity-driven RL to incentivize pixel-space exploration.

While these methods share \algoname's core principle of treating perceptual input as an active resource to be iteratively queried rather than as static context, time series reasoning differs from visual search in both problem setting and learning signal. Vision tasks typically target entities with well-defined spatial boundaries and often benefit from explicit supervision such as bounding boxes or segmentation masks, whereas temporal segments lack task-agnostic annotations and their relevance is inherently question-dependent, requiring \algoname to learn selection through weak supervision from downstream reasoning outcomes. Segments are also not self-contained: a zoomed image region is often interpretable in isolation, whereas a temporal segment's meaning depends on other segments - a spike is informative only relative to a baseline, and a regime shift requires comparing intervals before and after, so selection must be sequential and context-aware. \algoname's controller-reasoner decomposition and reliability-based reward target these challenges directly.


\xhdr{Self-Play RL for Reasoning}
Self-play originated in game theory and reinforcement learning as a mechanism for learning in adversarial multi-agent settings, where agents improve by interacting with copies of themselves or past versions of their policies rather than with a fixed opponent, and where outcomes depend directly on the evolving strategies of other agents, requiring continual adaptation instead of optimization against a static objective ~\cite{selfplay_survey}.
More recently, this paradigm has been adapted to LLMs, where self-play is used to generate training signals and improve reasoning without relying on human-labeled data~\cite{zhao2505absolute, fang2025serl,liu2025spice,yang2025spell, huang2025r, chen2025self, he2025visplay, wang2025vision, yu2025guided, yue2026dr, wang2026v}.  In these studies, a single LLM assumes multiple roles within the same training loop, typically forming a difficulty-driven interaction. A \emph{proposer} generates tasks that challenge the model’s current capabilities, while a \emph{solver} attempts to solve them and is rewarded based on correctness. This alternation exploits the stateless nature of LLMs, allowing the same underlying model to perform distinct roles via prompting, thereby enabling data- and memory-efficient training without maintaining multiple separate models. As a result, self-play provides a  mechanism for refining model behavior through role-based interaction and outcome evaluation, rather than reliance on external human annotations. 
This is important for time-series reasoning and adaptive segment selection, where annotations specifying which temporal segments are relevant to a given question are difficult to obtain. \vspace{0.3em}

In contrast to adversarial self-play, \algoname uses collaborative self-play to separate time-series reasoning into distinct roles. \algoname assigns adaptive segment selection to a controller and answer construction to a reasoner. The controller selects temporal segments. The reasoner produces intermediate reasoning steps and the final answer conditioned on the selected segments. This separation distinguishes evidence acquisition from answer generation and supports role-specific learning signals.
%
Existing self-play methods, including AZR~\cite{zhao2505absolute}, SPICE~\cite{liu2025spice}, and SPELL~\cite{yang2025spell}, optimize each role with myopic objectives that reward the best immediate response within a single round. These objectives do not match temporal segment selection, where the model must build an informative set of segments across multiple rounds. \algoname addresses this mismatch with hierarchical optimization. The controller receives trajectory-level supervision with credit assigned across interaction rounds, while the reasoner is optimized at the final round conditioned on the selected segments.
